\begin{ledgroup}
\footnotesize
\pstart
\noindent
In den nachfolgenden zwei Stücken, der Aufzeichnung N.\;\ref{60030} und dem Konzept N.\;\ref{60100},
befasst sich Leibniz wieder mit den kontroversen Thesen des Jesuiten
\indexN{\textls{Vanni}, Giovanni Francesco 1638\textendash1709}Giovanni Francesco Vannis
und der von ihm aufgeworfenen Frage der korrekten Zerlegung des Gewichts einer Kugel auf zwei schiefen Ebenen (\textsc{G.\,F.\;Vanni}, \cite{01995}\textit{Specimen libri de momentis gravium}, Rom 1684 und  \cite{02000}\textit{Exegeses physico-mathematicae}, Rom 1685). 
Diese hatten Leibniz bereits zu mehreren Publikationen und Konzepten veranlasst (siehe dazu die Stücke N.\;\ref{intro_vanni-specimen}, N.\;\ref{intro_vanni-demonstratio}, \ref{41199}, \ref{intro_vanni-solutio-problematis} und \ref{intro_vanni-investigatio}).
In der Aufzeichnung N.\;\ref{60030} führt Leibniz zwei Theoreme ein, von denen er nur das (allgemeinere) "Theorema I" tatsächlich beweist. Diese diente ihm als Vorlage für das weit umfangreichere Konzept N.\;\ref{60100}, worin er beide Theoreme aus N.\;\ref{60030}, allerdings in umgekehrter Benennung und Reihenfolge, ausführlich behandelt.
Daraus ergibt sich zunächst die relative Datierung der zwei Stücke: Erst wurde N.\;\ref{60030}, dann N.\;\ref{60100} verfasst.
\pend
\pstart
Die vorgeschlagene Datierung beruht auf dem Wasserzeichen und Gegenmarke im Träger von N.\;\ref{60100}. Dieselben Zeichen sind auch  in einem Konzept mit dem Titel "Considerationes de vi rebus insita" (LH XXXV 4, 5, 11 Bl.\;15--16) und einer verbesserten Reinschrift desselben, ebd. Bl.\;17--18, belegt. Die Texte, zwei Vorarbeiten zu Leibnizens Aufsatz \cite{02063}\textit{De ipsa natura sive de vi insita actionibusque Creaturarum, pro Dynamicis suis confirmandis illustrandisque} (\cite{01023}\textit{AE} September 1698, S.\;427--440), sind nach heutigem Kenntnisstand auf die Jahre 1697/1698 datierbar.
Mangels weiterer Belege für das Wasserzeichen im Leibnizkorpus erscheint die Annahme einer Verwendungsspanne von einigen Jahren angemessen, weshalb sich die vorgeschlagene, wenngleich nur wahrscheinliche Datierung 1695 bis 1700 ergibt.
\pend

\pstart
Im Einklang mit dieser Einschätzung sind auch Leibnizens Angaben im Konzept N.\;\ref{60100}:
Er bezieht sich ausdrücklich auf die Thesen, die \indexN{\textls{Vanni}, Giovanni Francesco 1638\textendash1709}G.\,F.~Vanni in seinem Buch "de Momentis Gravium" (d.\,h. \textit{Exegeses physico-mathematicae, de momentis gravium, de vecte, ac de motu aequabiliter accelerato} von 1685) "ante annos complures" aufgestellt hatte (\refnote{FnC:LH_35_10_13_001-002-complures}).
Die Schwierigkeit ("difficultas"), die Vanni als Gegenbeweis der bekannten Grundsätze der Statik angeführte Beispiel der Kugel auf zwei schiefen Ebenen anführte, habe Leibniz zwar bereits lösen können ("Ego memini me olim soluisse difficultatem"), aber die alte Schrift ("vetus scheda") sei ihm derzeit nicht zur Hand (\refnote{FnC:35_10_13_001-002-scheda}).
Dabei handelt es sich um die unveröffentlicht gebliebenen Konzepte \textit{Solutio problematis R.\,P.\;Vannii S.\,J.} von April bis Juni 1690 (N.\;\ref{41823}), die gänzlich der Lösung des Problems geweiht sind.
Die Angaben "ante annos complures" für 1685 und "vetus scheda" für die \textit{Solutio problematis} von 1690 stimmen mit der vorgeschlagenen Datierung der vorliegenden Stücke zwischen 1695 und 1700 überein.
\pend
\end{ledgroup}
